% SPDX-License-Identifier: CC-BY-4.0
% Copyright (C) 2026 Fabian Januszewski, Abhijit (the Langlandscape team)
% GENERATED by tools/extract_statements.py from mathematics/chapters/ -- do not edit by hand.

\begin{quoted}{Lemma}{13.2}{109}{lem:s04-ekills}{$e$ kills degree, vertex and component terms; seed Lemma ekills}{chapters/s04/ordinary.tex}
Let $s\ge1$, $V=\Zp$, $U$ of level prime to $p$.
\begin{enumerate}[label=(\alph*)]
\item $\deg:\Zp[X_s(U)]\to\Zp$ satisfies $\deg\circ U_p=p\cdot\deg$. Hence $e$ kills $\Zp$ (with $U_p=p$), $e\,\Zp[X_s]^0=e\,\Zp[X_s]$,
  and, for $q\nmid Dp$ and $U_q=\GL_2(\Zq)$, $e$ kills $H_0(\Gamma_{UU_s},\Zp)\cong\Zp$ (with the $U_p$-action induced from
  $\coker\partial$, \texttt{if:s02-coh}).
  \provisional{S04-M1: the $H_0$ clause uses that $\BG(UU_s,q)$ is connected (as in Thm.~lambdaribet, F7); the stage tex's proof
  says ``the connected graph'' but the statement does not list the hypothesis.} Mathematically the hypothesis is a theorem:
  $\Nrd(UU_s)=\hat\Z^\times$ and $B_q$ is split, so strong approximation for $B^1$ (Eichler, Kneser) gives connectivity; it is kept as
  a hypothesis because the Lean side carries it as \texttt{Hyp.SA}. (Also, $H_0(\Gamma,\Zp)\cong\Zp$ holds for any group; connectivity
  is needed for its identification with $\coker\partial$.)
\item ($B=B^{(q)}$, the definite algebra ramified exactly at $\{q,\infty\}$, written $B_q$ in stage~03, Def.~8.6; not
  the local algebra $B\otimes\Q_q$ of (N4); here $U=\Uone M$, so that $UU_s=\Uone{Mp^s}$ is a level of Def.~8.6, defined on
  $\hat B^{(q)\times}$ by the same local conditions; DR graph $\GDR_q(UU_s)$, \pathtt{if:s03-graph}.) Let $U_p$ act on the vertex module $\Zp[\{v_0,v_1\}]$ by $p$; then
  $\partial$ is $U_p$-equivariant, $e$ kills the vertex term, and $e\,\Xcomb_q(UU_s)\otimes\Zp=e\,\Zp[\XBq{UU_s}]$.
\item ($B=B^{(q)}$ as in (b).) $\Phicomb(\GDR_q(UU_s))\otimes\Zp\cong\Zp/|X_s|$ carries the action $\varphi\mapsto\varphi\circ U_p^*$ induced from
  $\Hom(X,\Z)$, which is multiplication by $p$; so $e(\Phicomb\otimes\Zp)=0$. This is a non-trivial statement: $v_p|X_s|\ge 2s-2$.
\item (Any $B$.) The Brandt pairing $\langle[x],[y]\rangle=\delta_{xy}$ on $\Zp[X_s]$ restricts to a perfect pairing
  \[e\,\Zp[X_s]\times e^*\Zp[X_s]\to\Zp,\]
  and $\nu$ identifies $e^*\Zp[X_s]$ with $e\,\Zp^{X_s}$ for the \emph{cohomological} $U_p$ ($f\mapsto\sum_jf(\cdot\,\varpi_j)$).
  Thus the perfect pairing is the evaluation pairing $e\,\Zp[X_s]\times e\,\Zp^{X_s}\to\Zp$. ($e$ and $e^*$ differ in general:
  Remark~13.3.)
\end{enumerate}
Here $U_p^*=T_{\varpi_p^{-1}}=\langle p\rangle_MU'_p$, where $U'_p=[UU_s\,\eta_p\,UU_s]$, $\eta_p=\iota_p^{-1}\mathrm{diag}(1,p)$, and
$\langle p\rangle_M$ is right translation by the central idele $t$ with $t_\ell=p$ for $\ell\ne p$ and $t_p=1$; for
$U^{(p)}=\Uone M$ this is the diamond $\langle d\rangle$ with $d\equiv p\bmod M$, $d\equiv1\bmod p^s$ (Lemma~4.3(d)).
\stproved{} below ((a)--(c) as in the seed, with the missing equivariance details; (d) with the identification of $e^*$; the
$H_0$ clause of (a) under the connectivity hypothesis). \emph{Lean} (for $\Zp$, $U=\Uone M$, $p\ge5$; the argument \texttt{hdec} is discharged by
\pathtt{padicOrdDecomp\_holds}): (a), degree part, \stformal: \lean{Langlandscape.Ordinary.deg\_comp\_UpHom},
\lean{Langlandscape.Ordinary.ekillsA\_tower} ($e\,\Zp[X_s]^0=e\,\Zp[X_s]$; abstract form \lean{Langlandscape.Ordinary.ordPart\_ker\_deg}); (d), perfect pairing,
\stformal: \lean{Langlandscape.Ordinary.perfect\_ordPart\_pairing\_tower\_padic}, \lean{Langlandscape.Ordinary.brandtPairing\_eq\_deltaPairing}. For (b), (c) only
the \emph{abstract} parts are formalized (\lean{Langlandscape.Ordinary.ordPart\_le\_ker\_of\_vertex}, \lean{Langlandscape.Ordinary.ordPart\_tensor\_phiComb\_eq\_bot}, with the
graph inputs and the identification $\Phicomb\cong\Z/|X_s|$ as hypotheses; the valuation bound $v_p|X_s|\ge2s-2$ modulo \pathtt{Hyp.EichlerMass} and a level-index input
\texttt{hlev}). Not formalized: the $H_0$ clause, the flatness identification $\Xcomb\otimes\Zp=\ker\partial_{\Zp}$, and the identification of $e^*$ with $e\,\Zp^{X_s}$
under $\nu$. So the lemma as a whole stays \stproved{} (tex).
\end{quoted}

\begin{quoted}{Theorem}{14.4}{114}{thm:s04-control}{$\Lambda$-freeness and control; seed Prop.~free (b); \texttt{if:s04-Mord}}{chapters/s04/control.tex}
Let $U$ be of level prime to $p$ and $r:=\rk_{\Zp}e\,\Zp[X_1(U)]$.
\begin{enumerate}[label=(\alph*)]
\item $\Mord(U)$ is a free $\Lam$-module of rank $r$, and for every $s\ge1$,
  $\Mord(U)\otimes_\Lam\cOE[\Delta_s]\xrightarrow{\sim}e\,\cOE[X_s(U)]$, compatibly with all Hecke operators.
\item (weights; $U=\Uone M$) For $k\ge2$, $\Mord(U)\otimes_\Lam\Lam/P^{\rm ad}_k$ is dual to $e\,\Brandt{U\Kp1}{\Vk\otimes\cOE}$ (Def.~16.1).
  On the component where $\mu_{p-1}$ acts by $\omega^{-(k-2)}$ it is dual to $e\,\Brandt{U\cdot\iota_p^{-1}\Uzero{p}_p}{\Vk\otimes\cOE}$, for every $k\ge2$.
  For $k\ge3$ the latter equals $e\,\Brandt{U}{\Vk\otimes\cOE}$ by $p$-stabilisation (\stexp{} in our normalisation).
  The first sentence is claimed integrally (over $\cOE$) only if $p\nmid\varphi(M)$; for $p\mid\varphi(M)$ it holds after inverting $p$,
  and its integral form is \stsketch{} (Remark~16.2(iv)). For a general level $U$ prime to $p$, (b) is \stexp{}
  (source: Hida 1988, Thm.~9.4, not checked here).
  \provisional{F4: (b) restated so that the character piece and the level match; before: ``For $k\ge2$ and $\chi$ the component on which $\mu_{p-1}$ acts by $\omega^{-(k-2)}$, $\Mord(U)^{(\chi)}\otimes_\Lam\Lam/P^{\rm ad}_k$ is dual to $e\,M(U\Kp1;\Vk\otimes\cOE)$; for $k\ge3$ one may replace $U\Kp1$ by $U\cdot\iota_p^{-1}\Uzero p_p$ (Hida's control at $\Gamma_0(p)$-level; expected)''; the left side was one $\mu_{p-1}$-eigenpiece and the right side the whole level-$U\Kp1$ module.}
\end{enumerate}
Status: (a) \stproved{} from Lemma~14.3 and Prop.~14.2 (choose compatible bases by Nakayama;
$\Lam=\varprojlim\cOE[\Delta_s]$); proof below, which uses only stage-01 results and elementary commutative algebra (raised from
\stpmc{} after the chapter audit S04-M9). Also \stcited{} \lit{hsieh\_2017\_hida-families-and-p-adic-triple-product-l-functions\_arxiv}{Thm.~4.2 and
its proof, p.~32: $\varprojlim_\alpha e\cO[X_\alpha]$ is $\Lam[\Delta_p]$-free, via [Hid88b, Cor.~10.1]} (Hsieh's ``Cor.~10.1'' is
\emph{Theorem}~10.1 in Hida's paper; there is no Cor.~10.1) and
\lit{longo-vigni\_2011\_control-theorems-for-quaternionic-hida-families\_arxiv}{Prop.~3.6 and its proof, p.~28: $\varprojlim_r
S_2(U_r,\cO_f)^{\rm ord}$ is $\Lam$-free by [Hid88b, Cor.~10.4]}.
(b) for $U=\Uone M$: \stcited{}, integrally if $p\nmid\varphi(M)$ (otherwise the integral form is \stsketch); general $U$: \stexp.
\lit{hsieh\_2017\_hida-families-and-p-adic-triple-product-l-functions\_arxiv}{Thm.~4.2(2), p.~32 (forms side; via
[Hid88b, Thm.~9.4])}; \lit{longo-vigni\_2011\_control-theorems-for-quaternionic-hida-families\_arxiv}{Thm.~3.7, p.~28 (eigen-components,
over a field)}; the dictionary to our $\NUp$, $\Vk$ and $P^{\rm ad}_k$ is Remark~16.2 (S04-M4), whose key
computation is Def.~16.1. The original source is
\lit{hida\_1988\_p-adic-hecke-algebras-for-gl2-over-totally-real-fields}{Thm.~8.6, p.~362 (weight independence of $\mathcal V^{\rm ord}$);
Thm.~9.4, p.~369 (control); Thm.~10.1, p.~370 (freeness over $\cO[[Z_a]]$ under $p$-divisibility); Cor.~10.4, p.~373 ($\Lam$-freeness,
$p>2$)}: the theorem headers and page numbers are located in the OCR text, but the statements are only partly legible, so their exact
hypotheses are taken from Hsieh and LV11.
\emph{Lean status of (a):} the finite-level part is \stformal{} for $\Zp$, $U=\Uone M$, $p\ge5$: control as surjectivity plus kernel
(\lean{Langlandscape.Ordinary.ord\_control\_finite}, \lean{Langlandscape.Ordinary.TowerHyp.mordTrans\_surjective},
\lean{Langlandscape.Ordinary.Mord.proj\_surjective}), freeness at level $s$ (\lean{Langlandscape.Ordinary.ordFree\_finite}). \emph{Not} formalized:
the constancy of the rank in $s$ (\pathtt{ordRank\_const}); the $\Lam$-adic half ($\Mord$ free of finite rank over $\Lam$ and the control
isomorphism for all $s$) is only stated, as the open hypothesis \lean{Langlandscape.Hyp.MordLambdaFree} (Lean-side status: open, S04-L14b). Thus
(a) as a whole stays \stproved{} (tex), with the Lean part \emph{formalized modulo \pathtt{Hyp.MordLambdaFree}} for the $\Lam$-adic half. The
Lean hypothesis \lean{Langlandscape.Hyp.MordLambdaFree} is thus the statement of a tex-proved theorem (not a merely cited one), taken as an argument because the data
$G_n=\Delta_{n+1}$, the reduction maps and the diamond representations are not yet in Lean; it has no named consumer in the stage-05/06 plans.
\emph{(b), Lean form (M9.3 sub-check, recorded):} \lean{Langlandscape.Hyp.HidaControl} is schematic, on abstract \emph{modules} $M_k$, $M_2$, $Q_k$ (the tex theorem is stated on
duals, $P^{\rm ad}_k$-quotients); it carries the guard $p\nmid\varphi(M)$ (M9.3) and $U_1(M)$-type levels only; its arguments \texttt{hk} and \texttt{h$\varphi$} are not used in the
statement (audit item O1 not applied).
\end{quoted}

\begin{quoted}{Theorem}{15.1}{116}{thm:s04-lambdaribet}{$\Lambda$-adic definite Ribet sequence; seed Thm.~lambdaribet; \pathtt{if:s04-lambdaribet}}{chapters/s04/lambdaribet.tex}
Let $q\nmid Dp$, $U$ of level prime to $pq$ with $U_q=\GL_2(\Zq)$ ($q$-admissible in the sense of stage 02). Assume $\BG(UU_s,q)$ is
connected for every $s\ge1$ (equivalently, Prop.~quot holds at all these levels). This follows from \texttt{Hyp.SA} $\wedge$
\texttt{Hyp.HS} (in Lean also \texttt{Hyp.BTtree}); it is not a finite check, and it does not imply $H_0=V_{U_p^1}$.
\provisional{F7: connectivity hypothesis made explicit; before: no such hypothesis was stated (only the implicit ``$q$-admissible'').}
Here $\Gamma_{UU_s}:=(\cO_{UU_s}[1/q])^1$ is the norm-one $\{q\}$-arithmetic group (as in Step~2 of the proof of Prop.~18.1); it does not
contain the central $q^{\Z}$. Put
\[\Xord_{Dq}(U):=\varprojlim_se\,H_1(\Gamma_{UU_s},\Zp)\otimes\cOE\]
(transition maps: corestriction, i.e.\ the restriction of $\pi_*$ on $\Zp[X^{(q)}_s]$-chains; over $\Zp$, $p\ge5$,
$H_1(\Gamma,\Zp)=H_1(\BG,\Z)\otimes\Zp$ by stage 02, D1). Then:
\begin{enumerate}[label=(\alph*)]
\item the sequence of $\Lam$-modules, equivariant for all Hecke operators supported away from $q$ (including $U_p$, $\langle d\rangle$,
  $\Pi_\ell$ for $\ell\mid D$),
  \[0\to\Xord_{Dq}(U)\to\Mord(\Uzero{q}U)\xrightarrow{\ \alpha=(-s_{0*},\,s_{1*})\ }\Mord(U)^{\oplus2}\to0\]
  is exact;
\item all three terms are $\Lam$-free, and $\rk_\Lam\Xord_{Dq}(U)=\rk_\Lam\Mord(\Uzero qU)-2\rk_\Lam\Mord(U)$;
\item at finite level $s\ge1$: $\rk_{\Zp}H_1(\Gamma_{UU_s},\Zp)=(q-1)|X_s(U)|+1$, and $\rk_{\Zp}e\,H_1=\rk e\Zp[X^{(q)}_{s}]-2\rk e\Zp[X_s]$,
  the ``$+1$'' being the degree line killed by $e$ (here $X^{(q)}_s$ is the level $\Uzero qUU_s$).
\end{enumerate}
Status: (a) \stproved{} below given \pathtt{if:s03-ribet}, Lemma~13.2(a) and $\NUp$; (b) \stproved{} via
Thm.~14.4(a); (c) \stproved{} (seed rank check; \pathtt{if:s02-indef}). \emph{Lean status:} \stformal{} for $E=\Q_p$, $U=\Uone M$, $p\ge5$, $q\ne p$:
the finite-level sequence of (a) and (c) at every $s\ge1$ (\lean{Langlandscape.Ordinary.ordRibet\_exact\_Utower\_uncond},
\lean{Langlandscape.Ordinary.rank\_ordRibet\_uncond}, \lean{Langlandscape.Ordinary.rank\_ordRibet\_full\_uncond}; connectivity of $\BG(UU_s,q)$ is an explicit
argument, \pathtt{Hyp.PadicOrdDecomp} is discharged). The limit statement of (a) is \emph{formalized modulo \pathtt{Hyp.InvLimExact}}
(and the connectivity argument): \lean{Langlandscape.Ordinary.lim\_exact\_ribet} proves injectivity, exactness in the middle and surjectivity of $\alpha$ on the
limits, but takes the surjectivity of the $A$-transitions $e\,H_1(\Gamma_{UU_{s+1}})\to e\,H_1(\Gamma_{UU_s})$ as the literal argument
\texttt{hA : $\forall n$, Function.Surjective (H.transA n)} (not a term of type \lean{Langlandscape.Hyp.InvLimExact}: that hypothesis is the implication
``$\forall s\ge1$ \texttt{BGConnected} $\Rightarrow\forall n$, \texttt{Surjective (H.transA n)}'', so \texttt{hA} is its conclusion, available once the connectivity argument is supplied). This $A$-surjectivity is \emph{not} formalized (undischarged; the snake-lemma route was dropped in S04-L13b), so
(a) stays \stproved{} (tex) and is not upgraded. Not formalized: (b) (the $\Lam$-rank formula).
\end{quoted}

\begin{quoted}{Theorem}{16.6}{122}{thm:s04-Gor}{Gorenstein and freeness; \texttt{if:s04-Mord}, second half}{chapters/s04/hsieh.tex}
Let $p>3$, $p\nmid MD$, $N=MD$, and let $F\in eS(N,\chi,\mathbb I)$ be an ordinary $\Lam$-adic newform of tame conductor $N$ with
$\mathrm{supp}\,N^-\subseteq\Sigma^0_F$ (discrete series at every $\ell\mid D$; Hsieh's hypothesis for the factorisation of
$\lambda_F$ through JL, [Hs, \S4.5, p.~34]), satisfying (CR, $\mathrm{supp}\,N^-$) (Hyp.~16.5). Let $\m_F\subset\T^D(N,\mathbb I)$
be the maximal ideal containing $\ker\lambda^D_F$ (above the maximal ideal $\m$ of Hyp.~16.5; all localisations below are
at $\m_F$, of $\T^D(N,\mathbb I)$-modules). Then:
\begin{enumerate}[label=(\roman*)]
\item Hsieh's $eS^D(N,\mathbb I)_{\m_F}$ is free of rank one over $\T^D(N,\mathbb I)_{\m_F}$ (proof of [Hs, Thm.~4.5]);
\item $eS^D(N,\chi,\mathbb I)_{\m_F}$ is free of rank one over $\T^D(N,\chi,\mathbb I)_{\m_F}$, and $\T^D(N,\chi,\mathbb I)_{\m_F}$ is
  Gorenstein (ibid.);
\item the eigenline $eS^D(N,\mathbb I)[\lambda^D_F]$ is free of rank one over $\mathbb I$ (its generator $F^D$: the primitive JL lift)
  ([Hs, Thm.~4.5]);
\item if moreover $p\nmid\varphi(M)$, then $\T^D(N,\mathbb I)_{\m_F}=\T^D(N,\chi,\mathbb I)_{\m_F}$ is Gorenstein, and
  $(\Mord(\Uone M)\otimes_\Lam\mathbb I)_{\m_F}$, which is the $\iota_2$-twisted $\mathbb I$-dual of $eS^D(N,\mathbb I)_{\m_F}$
  (Remark~16.3), is free of rank one over $\T^D(N,\mathbb I)_{\m_F}$.
\end{enumerate}
If $p\mid\varphi(M)$, Gorenstein-ness of the full $\T^D(N,\mathbb I)_{\m_F}$ and the homological freeness in (iv) are not cited and
not proved here (\stexp). For other $\m$ satisfying (CR$_D$) (old at $M$, or of smaller conductor) the conclusions are \stexp{}
(cf.\ [Wil95, Thm.~2.1], [Hel07, Cor.~8.11]; not checked here).
\par\smallskip\noindent{\scriptsize\itshape [Status paragraph following the statement in the chapter:]}\par
\noindent Registered in Lean as the hypothesis \lean{Langlandscape.Hyp.HsiehGor} (cited, not discharged; used by \lean{Langlandscape.Defect.Setting.gor\_new}, stage~6, which instantiates it over $\Lambda$).
\noindent\stcited{} (i)--(iii) \lit{hsieh\_2017\_hida-families-and-p-adic-triple-product-l-functions\_arxiv}{Thm.~4.5 and its proof, p.~34};
(iv) \stproved{} from (ii): if $p\nmid\varphi(M)$, the torsion of the diamond group, of order $(p-1)\varphi(M)$, is prime to $p$, so
(after enlarging $\cO$) its idempotents lie in $\T^D(N,\mathbb I)$ and $eS^D(N,\mathbb I)=\bigoplus_{\chi'}eS^D(N,\chi',\mathbb I)$; the
characters $\chi'$ have distinct reductions (orders prime to $p$), so exactly one idempotent, that of $\chi$, survives in the local ring
$\T^D(N,\mathbb I)_{\m_F}$, giving $\T^D(N,\mathbb I)_{\m_F}=\T^D(N,\chi,\mathbb I)_{\m_F}$; dualising a free rank-one module over a
Gorenstein ring gives a free rank-one module. The proof of Thm.~4.5 cites [Wil95, Thm.~2.1], [Hel07, Cor.~8.11, Rem.~8.12]
(\lit{helm\_2004\_maps-between-modular-jacobians-and-jacobians-of-shimura-curves\_arxiv}{Cor.~8.11 and Rem.~8.12, p.~31; journal
version p.~110, not checked on this node}) and [PW11, Prop.~6.4, 6.5]
(\lit{pollack-weston\_2006\_anticyclotomic-mu-invariants-of-modular-forms\_arxiv}{Prop.~6.4, p.~19; Prop.~6.5, p.~20}); the
freeness statement itself there is PW Thm.~6.2, p.~19, transported to the definite side by Prop.~6.5 (Prop.~6.4 is the self-pairing
statement, conditional on freeness). The inputs carry their own hypotheses (PW's CR; Helm's ``controllable'' $\m$; Wiles' Thm.~2.1);
that (CR, $\mathrm{supp}\,N^-$) suffices for this chain is Hsieh's claim and is not re-verified here. Gorenstein
follows from freeness plus the perfect pairing $B_N$. The freeness in (i), (ii) is in the \emph{proof}, not the statement, of
Thm.~4.5 (seed cites only ``Thm.~4.5''), and Hsieh states Gorenstein only for the $\chi$-piece $\T^D(N,\chi,\mathbb I)_{\m_F}$. Passing between $\Mord$ (homological) and $eS^D$ (forms) uses Hsieh (4.6) and
Thm.~4.2(1); dualising a free rank-one module over a Gorenstein ring is again free of rank one.
\end{quoted}

\begin{quoted}{Proposition}{17.1}{124}{prop:s04-etaprinciple}{$\eta$-principle; \pathtt{thm:s04-etaB}, \texttt{lem:s06-S}}{chapters/s04/etaprinciple.tex}
Let $A$ be a commutative ring and $S$ a commutative $A$-algebra that is finite free as an $A$-module. Let $M$ be an $S$-module with an
$S$-linear isomorphism $e:M\xrightarrow{\sim}S$, and let $B:M\times M\to A$ be an $A$-bilinear form such that
\begin{itemize}
\item $B$ is \emph{perfect}: the $A$-linear map $M\to\Hom_A(M,A)$, $m\mapsto B(m,\cdot)$, is bijective;
\item $B$ is \emph{$S$-equivariant}: $B(sm,n)=B(m,sn)$ for all $s\in S$, $m,n\in M$.
\end{itemize}
Let $\mu:S\to A$ be a homomorphism of $A$-algebras, and let $m_0\in M$ generate the eigenline
\[M[\mu]:=\{m\in M:\ sm=\mu(s)m\ \text{for all }s\in S\},\qquad\text{i.e.\ }M[\mu]=A\,m_0.\]
Then:
\begin{enumerate}[label=(\alph*)]
\item $S$ is Gorenstein with trace $\ell(s):=B(m_1,s\,m_1)$, where $m_1:=e^{-1}(1)$ (\lean{Langlandscape.CommAlg.EtaPrinciple.isGorensteinWith\_of\_pairing});
\item $\mathrm{congruenceIdeal}(\mu)=\big(B(m_0,m_0)\big)$, the principal ideal of $A$ generated by $B(m_0,m_0)$
  (\lean{Langlandscape.CommAlg.EtaPrinciple.congruenceIdeal\_eq\_span\_pairing}).
\end{enumerate}
\stproved{} below. \stformal{} (S06-L17, \pathtt{CommAlg/EtaPrinciple.lean}), in full generality of the statement and without registered hypotheses: \lean{Langlandscape.CommAlg.EtaPrinciple.isGorensteinWith\_of\_pairing} for (a) and \lean{Langlandscape.CommAlg.EtaPrinciple.congruenceIdeal\_eq\_span\_pairing} for (b); the trace in Lean is \texttt{traceOfPairing}, $\ell=B(e^{-1}(1),e^{-1}(\cdot))$.
\end{quoted}
